\documentclass[11pt]{article}
\usepackage{amsmath}
\usepackage{amssymb}
\usepackage[margin=2cm]{geometry}
\usepackage{longtable}
\usepackage{booktabs}
\usepackage{pdflscape}
\usepackage{breqn}
\title{Certificate for [pi]\_q (first 4 coefficients)}
\author{qreals certificate}
\date{\today}
\begin{document}
\maketitle

\noindent Input: the real number x = pi. The continued fraction is truncated to the depth that locks the first 4 coefficients (MGO Proposition 1.1).

\section*{(a) Continued fraction and even-length MGO form}
Regular continued fraction: [3, 7]. \par
Even-length MGO form: [3, 7]. \par
The even-length form evaluates to the convergent $ \frac{22}{7} $.\par

\section*{(b) MGO formula folded step by step}
Odd positions carry $[a]_q$ with $q^a$ above; even positions carry
$[a]_{q^{-1}}$ with $q^{-a}$ above (MGO eqn.~1.1).

\noindent\textbf{innermost term a\_2 = 7:}
{\footnotesize%
\begin{dmath*}
\frac{q^{6} + q^{5} + q^{4} + q^{3} + q^{2} + q + 1}{q^{6}}
\end{dmath*}
}

\noindent\textbf{fold in a\_1 = 3:}
{\footnotesize%
\begin{dmath*}
\frac{q^{9} + q^{8} + 2 q^{7} + 3 q^{6} + 3 q^{5} + 3 q^{4} + 3 q^{3} + 3 q^{2} + 2 q + 1}{q^{6} + q^{5} + q^{4} + q^{3} + q^{2} + q + 1}
\end{dmath*}
}

\noindent\textbf{Result.} [22/7]\_q (the convergent, whose Taylor expansion gives [x]\_q):
{\footnotesize%
\begin{dmath*}
\frac{q^{9} + q^{8} + 2 q^{7} + 3 q^{6} + 3 q^{5} + 3 q^{4} + 3 q^{3} + 3 q^{2} + 2 q + 1}{q^{6} + q^{5} + q^{4} + q^{3} + q^{2} + q + 1}
\end{dmath*}
}

\noindent Taylor coefficients:
{\small%
\begin{dmath*}
1 + q + q^{2} + O(q^{4})
\end{dmath*}
}

\section*{(c) Independent cross-checks}
\begin{itemize}
\item \textbf{[pass]} the first 4 coefficients are unchanged when 8 are asked for
\item \textbf{[pass]} [x+1]\_q computed from its own continued fraction equals q*[x]\_q + 1
\item \textbf{[na]} the series [pi]\_q diverges at q=1; q=1 is for rationals
\item \textbf{[na]} pi is irrational, so the continued fraction does not terminate
\end{itemize}
Summary: verified: truncation stable to 4, shift law [x+1]=q[x]+1; n/a: q=1 specialisation; exact rational function.

\section*{References}
S. Morier-Genoud and V. Ovsienko, q-deformed rationals and q-continued
fractions, Forum Math. Sigma \textbf{8} (2020), e13. Per-function theorem and
check mapping: \texttt{docs/CORRECTNESS.md}.

\medskip
\noindent This is a human-auditable derivation, not a formal machine proof;
every line above is checkable by hand.
\end{document}
